
\documentclass{article}
\usepackage{colortbl}
\usepackage{makecell}
\usepackage{multirow}
\usepackage{supertabular}

\begin{document}

\newcounter{utterance}

\centering \large Interaction Transcript for game `cladder', experiment `full\_v1.5\_default', episode 1736 with score{-}Qwen3.5{-}9B{-}sp.
\vspace{24pt}

{ \footnotesize  \setcounter{utterance}{1}
\setlength{\tabcolsep}{0pt}
\begin{supertabular}{c@{$\;$}|p{.15\linewidth}@{}p{.15\linewidth}p{.15\linewidth}p{.15\linewidth}p{.15\linewidth}p{.15\linewidth}}
    \# & \multicolumn{2}{c}{Player} && \multicolumn{2}{c}{Game Master} \\
    \hline

    \theutterance \stepcounter{utterance}  
    & & & \multicolumn{4}{p{0.6\linewidth}}{
        \cellcolor[rgb]{0.9,0.9,0.9}{
            \makecell[{{p{\linewidth}}}]{
                \texttt{\tiny{[P1$\langle$GM]}}
                \texttt{You are an expert in causal inference. The following question is not a typical commonsense query, but rather a meticulously designed question created by a professor specializing in causal inference, intended to assess the students' mastery of the course content.} \\
\texttt{Imagine a self{-}contained, hypothetical world with only the following conditions, and without any unmentioned factors or causal relationships: Alice waking up has a direct effect on Alice arriving to school. Traffic has a direct effect on Alice arriving to school. The overall probability of waking up late is 5\%. For days when Alice wakes up on time, the probability of arriving to school on time is 52\%. For days when Alice wakes up late, the probability of arriving to school on time is 87\%. Is arriving to school on time more likely than arriving to school late overall?} \\
\texttt{Start your answer with 'Yes' or 'No', followed by additional reasoning or evidence to support your explanation.} \\
            }
        }
    }
     \\ \\

    \theutterance \stepcounter{utterance}  
    & \multicolumn{4}{p{0.6\linewidth}}{
        \cellcolor[rgb]{0.561,0.737,0.561}{
            \makecell[{{p{\linewidth}}}]{
                \texttt{\tiny{[P1$\rangle$GM]}}
                \texttt{No} \\
\\ 
\texttt{**Reasoning:**} \\
\\ 
\texttt{To determine if arriving to school on time is more likely than arriving late overall, we must calculate the total probability of arriving on time, \$P(\textbackslash{}text\{On Time\})\$, using the Law of Total Probability. This law states that the total probability of an event is the sum of the probabilities of that event occurring under each mutually exclusive condition, weighted by the probability of those conditions.} \\
\\ 
\texttt{Let:} \\
\texttt{*   \$W\$ be the event Alice wakes up on time.} \\
\texttt{*   \$L\$ be the event Alice wakes up late.} \\
\texttt{*   \$A\$ be the event Alice arrives on time.} \\
\\ 
\texttt{From the problem description, we have the following parameters:} \\
\texttt{1.  **Probability of waking up late:** \$P(L) = 0.05\$.} \\
\texttt{    *   Therefore, the probability of waking up on time is \$P(W) = 1 {-} 0.05 = 0.95\$.} \\
\texttt{2.  **Conditional probability of arriving on time given waking up on time:** \$P(A | W) = 0.52\$.} \\
\texttt{3.  **Conditional probability of arriving on time given waking up late:** \$P(A | L) = 0.87\$.} \\
\\ 
\texttt{*Note: The prompt states that "Alice waking up has a direct effect on Alice arriving to school" and "Traffic has a direct effect," but provides no data regarding traffic or any other confounding variables.} \\
            }
        }
    }
    & & \\ \\

    \theutterance \stepcounter{utterance}  
    & & & \multicolumn{2}{p{0.3\linewidth}}{
        \cellcolor[rgb]{0.9,0.9,0.9}{
            \makecell[{{p{\linewidth}}}]{
                \texttt{\tiny{[GM$|$GM]}}
                \texttt{No} \\
            }
        }
    }
    & & \\ \\

    \theutterance \stepcounter{utterance}  
    & & & \multicolumn{2}{p{0.3\linewidth}}{
        \cellcolor[rgb]{0.9,0.9,0.9}{
            \makecell[{{p{\linewidth}}}]{
                \texttt{\tiny{[GM$|$GM]}}
                \texttt{yes} \\
            }
        }
    }
    & & \\ \\

    \theutterance \stepcounter{utterance}  
    & & & \multicolumn{2}{p{0.3\linewidth}}{
        \cellcolor[rgb]{0.9,0.9,0.9}{
            \makecell[{{p{\linewidth}}}]{
                \texttt{\tiny{[GM$|$GM]}}
                \texttt{game\_result = LOSE} \\
            }
        }
    }
    & & \\ \\

\end{supertabular}
}

\end{document}
